% paper/section_3_design.tex
% ─────────────────────────────────────────────────────────────────────────────
% Section 3: Kavach Architecture
% Slot into skeleton.tex where \section{Kavach Architecture} appears.
% Numbers marked [TBD] are filled in from benchmark workstream D output.
% ─────────────────────────────────────────────────────────────────────────────

\section{Kavach Architecture}
\label{sec:design}

\subsection{Overview}

Kavach is a runtime guardrail that interposes between an LLM agent orchestrator
and its tool-execution layer. Every tool call proposed by the agent is
intercepted, evaluated by a Parliament of domain-bound Ministers, and sentenced
by a deterministic Speaker before any side-effecting action is permitted to
execute. Figure~\ref{fig:architecture} shows the end-to-end request path.

The system exposes a single HTTP endpoint \texttt{POST /parliament/evaluate}
that the OpenClaw plugin (\S\ref{sec:plugin}) calls synchronously in the
\texttt{pre\_tool\_call} hook. The endpoint returns a JSON verdict
\texttt{\{decision, confidence, ministers\_fired, compass\_score,
drift\_delta\}} within the configured latency budget.

\subsection{Core Components}

\paragraph{Ledger.}
An append-only SQLite journal (\texttt{kavach\_parliament.db}) records every
evaluation request and its verdict. Each row captures:
\texttt{session\_id}, \texttt{turn}, \texttt{tool\_name}, \texttt{action\_embedding}
(1024-d float blob), \texttt{verdict}, and \texttt{timestamp\_utc}.
The Ledger is the single source of truth for session state; no in-memory
session object exists, making the server horizontally stateless between
requests. Session risk score is recomputed from the Ledger at each call via
a lightweight SQL aggregate, adding $<$2\,ms overhead.

\paragraph{Semantic Router.}
On each evaluation request the proposed action text is embedded with
\texttt{BAAI/bge-large-en-v1.5} (1024-d, cosine similarity). The resulting
vector is compared against five pre-computed domain centroids stored in
\texttt{kavach\_router\_config.json}. Ministers whose centroid cosine similarity
exceeds $\theta_{\text{route}} = 0.40$ are activated; COMPASS always runs
regardless of routing similarity.

The router replaces the stage-bound dispatcher used in early prototypes
(OpenClaw stages Init$\to$Input$\to$Inference$\to$Decision$\to$Execution).
Ministers are \emph{domain-bound}, not stage-bound: EXECUTOR fires on any
action resembling code execution regardless of which pipeline stage proposed it.
This is the key architectural departure from OpenClaw's native trust model
and from contemporaneous work such as NemoClaw~\cite{nvidia2025nemoclaw}, which
enforces a single boundary at the pipeline entry point.

\paragraph{Ministers.}
Each activated Minister runs a two-tier scan:

\begin{enumerate}
  \item \textbf{Syntactic tier (L1).} Regex and exact-match rules over the
        raw action string. Matches produce an \texttt{L1\_hit} flag and a
        severity label. High-severity L1 hits short-circuit to \textsc{Block}
        without invoking the semantic tier.
  \item \textbf{Semantic tier (L2/L3).} The action embedding is compared via
        cosine distance against every pattern in the Minister's ChromaDB
        collection. A hit is declared when
        $\cos(\mathbf{a}, \mathbf{p}) \geq \theta_{\min}$ where
        $\theta_{\min} = 0.45$ (Config\,B; see \S\ref{sec:eval}).
        The hit is labelled with the matching pattern's L1/L2/L3 descriptors
        for the verdict log.
\end{enumerate}

The four domain-bound Ministers and their attack-surface coverage:

\begin{itemize}
  \item \textbf{EXECUTOR} — code execution, process injection, persistence
        mechanisms, container escape (100 patterns; MITRE ATT\&CK T1059,
        T1055, T1053, T1611 and sub-techniques).
  \item \textbf{VAULT} — credential access, secret file reads, metadata service
        queries, platform keychain enumeration (100 patterns; T1552, T1528,
        T1555 and sub-techniques).
  \item \textbf{CHANNEL} — data exfiltration via any outbound channel: HTTP/S,
        DNS, SMTP, cloud storage, git push, covert channels (100 patterns;
        T1048, T1567, T1102 and sub-techniques).
  \item \textbf{NAVIGATOR} — session trajectory anomalies: post-task
        persistence, scope widening, recon sequences, anti-forensic actions
        (100 patterns; OWASP Agentic Top 10 2026 A01, T1083, T1082, T1070).
\end{itemize}

\paragraph{COMPASS.}
The Calibrated Oracle for Mission-Pattern Alignment and Session Semantics
evaluates \emph{intent alignment} rather than attack-pattern similarity. It
maintains a \emph{session centroid} — an exponential moving average of action
embeddings weighted by turn recency — and measures cosine drift between the
centroid and the seeded user-intent vector (extracted from the first
non-system turn of the session).

\begin{equation}
  \delta_{\text{drift}}(t) =
    1 - \cos\!\left(\bar{\mathbf{c}}_t,\, \mathbf{u}\right)
  \label{eq:drift}
\end{equation}

where $\bar{\mathbf{c}}_t$ is the turn-$t$ session centroid and $\mathbf{u}$ is
the user-intent embedding. When $\delta_{\text{drift}} > \theta_{\text{drift}}$
the call is escalated regardless of individual minister verdicts.

The COMPASS calibration set (\texttt{kavach\_corpus\_v1.json}, key
\texttt{COMPASS}) contains 100 labelled pairs — 50 \textsc{aligned} and 50
\textsc{hijacked} — covering 25 distinct user-intent scenarios each paired with
both a legitimate and an attack-injected proposed action. Youden's $J$
statistic is used to select $\theta_{\text{drift}}$ via
\texttt{compass\_calibrator.py}.

\paragraph{Speaker.}
The Speaker is a deterministic finite function — no learned parameters — that
maps the set of Minister verdicts and COMPASS output to a single sentence:

\begin{equation}
V = \text{Speaker}(\{v_m\}_{m \in M_{\text{active}}},\, \delta_{\text{drift}},\, r_{\text{session}})
\end{equation}

The decision table (five cases, applied in priority order):
\begin{enumerate}
  \item Any Minister returns \textsc{Block} $\Rightarrow$ \textsc{Block}.
  \item $\delta_{\text{drift}} > \theta_{\text{drift}}$ $\Rightarrow$ \textsc{Escalate}.
  \item Session risk $r_{\text{session}} > r_{\text{max}}$ $\Rightarrow$ \textsc{Escalate}.
  \item Any Minister returns \textsc{Grey} $\Rightarrow$ \textsc{Escalate}.
  \item All Ministers return \textsc{Allow} $\Rightarrow$ \textsc{Allow}.
\end{enumerate}

\subsection{OpenClaw Plugin}
\label{sec:plugin}

The Kavach plugin (\texttt{openclaw-plugin-kavach.ts}) implements the
OpenClaw plugin interface:

\begin{itemize}
  \item \textbf{\texttt{pre\_tool\_call}}: synchronously calls
        \texttt{/parliament/evaluate} before the agent executes a tool.
        On \textsc{Block} or \textsc{Escalate}, the plugin throws
        \texttt{KavachBlockError} or \texttt{KavachEscalateError}, preventing
        tool execution.
  \item \textbf{\texttt{on\_tool\_error}}: increments the circuit-breaker
        failure counter. After \texttt{kavachCircuitBreakerThreshold}
        consecutive Kavach server errors the plugin enters \emph{open-circuit}
        mode: all tool calls are blocked and an operator alert is raised.
\end{itemize}

The plugin is configured via \texttt{openclaw.plugin.json} and requires only
\texttt{KAVACH\_URL} (default: \texttt{http://localhost:8000}) and
\texttt{SESSION\_ID} to be present in the OpenClaw runtime environment.

\subsection{Design Decisions and Trade-offs}

\paragraph{Ministers are domain-bound, not stage-bound.}
OpenClaw's five-stage pipeline (Init$\to$Execution) provides natural hooks but
creates stage-specific blind spots: an attack that appears in the Inference
stage looks different from the same attack in the Decision stage, requiring
per-stage classifiers. By binding Ministers to attack \emph{domains} we achieve
consistent coverage across all stages with a single collection per domain.

\paragraph{Append-only Ledger, no session objects.}
Storing session state in an append-only log rather than an in-memory object
prevents an attacker from clearing session state via a crafted tool call.
The immutability guarantee is enforced at the database level
(\texttt{kavach\_parliament.db} opened read-write only by the server; no
\texttt{DELETE} or \texttt{UPDATE} statements are issued).

\paragraph{Synchronous evaluation on the hot path.}
Kavach evaluates in the \texttt{pre\_tool\_call} hook synchronously.
Asynchronous evaluation would reduce latency but allow a race between the
verdict and execution; we accept the latency cost to maintain the security
guarantee that \emph{no tool executes before a verdict is issued}.
